\section{结论}
\label{sec:conclusion}

本文对面向非专业用户的两种Vibe Coding策略进行了对比研究。
通过构建个人主页的自我实验，我们证明：增量构建——先搭建页面骨架再逐步添加模块——在早期即因缺乏统一的设计理念而失败。
没有设计锚点，每次Prompt都在独立运作，风格一致性无法维系，用户对最终产出感到失控。

相比之下，设计优先的Vibe Coding——在编码前先与AI共同设计完整规范——成功产出了一个具备复杂交互的完整多区域网站。
通过将设计构思与代码生成相分离，该策略将用户的角色从设计者转变为设计评审者，使非专业用户能够借助AI的设计能力，而非受限于自身的设计水平。

本文的核心发现如下：

\begin{enumerate}
    \item \textbf{设计构思不是额外开销，而是根基。}在编码前投入的20\%对话轮次用于设计讨论，防止了导致增量构建失败的风格漂移和沟通成本攀升。
    \item \textbf{将AI视为设计伙伴，而非代码机器。}最有效的Prompt是让AI\textit{提出设计方案}，而非直接生成代码。这种模式——AI提案、用户判断、AI优化——比用户直接指定设计参数更加高效。
    \item \textbf{设计文档具备双重用途。}它既是人类可读的规范（供用户验证设计意图），也是AI可读的约束（供模型在实现和纠错时参照）。
    \item \textbf{浏览器测试揭示设计文档无法捕捉的问题。}即使规范详尽，某些品质——视觉权重、动画手感、交互可发现性——只有通过实际使用才能显现。体验驱动迭代循环（测试$\rightarrow$发现偏差$\rightarrow$Prompt纠正$\rightarrow$验证）是设计优先策略的必要补充。
\end{enumerate}

本研究存在一定局限，包括单被试设计（一名用户、一个任务）以及策略与AI工具的偶发（增量构建使用Codex，设计优先使用Claude Code）。
未来工作可以控制工具变量的影响，引入多名不同设计能力的用户参与，并在静态网站之外的更复杂应用类型上测试设计优先策略。

对于实践者，我们的建议清晰明确：\textit{在让AI写代码之前，先让AI帮助你做设计。}
设计构思的投入通过降低修改开销和大幅提升产出质量来回报自身。
